\section*{Bab \ref{ch:g9:maps-of-the-earth} --- Peta Seluruh Bumi}

\begin{solution}{exo:g9:maps-of-the-earth:1}
(a) \omterm{def:g9:maps-of-the-earth:projection}{Proyeksi peta}. (b) \omterm{def:g9:maps-of-the-earth:contour}{Garis kontur}. (c) \omterm{def:g9:maps-of-the-earth:numerical-scale}{Skala angka}. (d) \omterm{def:g9:maps-of-the-earth:contour}{Interval
kontur}. (e) \omterm{def:g9:maps-of-the-earth:remote-sensing}{Penginderaan jauh}.
\end{solution}

\begin{solution}{exo:g9:maps-of-the-earth:2}
(a) $4 \times 25\,000 = \num{100000}$~cm $= \qty{1}{km}$. (b) $3 \times
50\,000 = \num{150000}$~cm $= \qty{1.5}{km}$. (c) $2.5 \times 1\,000\,000
= \num{2500000}$~cm $= \qty{25}{km}$. (d) $\qty{1}{km} = \num{100000}$~cm;
$100\,000 \div 25\,000 = \qty{4}{cm}$.
\end{solution}

\begin{solution}{exo:g9:maps-of-the-earth:3}
(a) Salah: permukaan bola tidak dapat diratakan tanpa mengubah
panjang-panjangnya, sehingga setiap \omterm{def:g3:continents-and-oceans:world-map}{peta dunia} mengubah sesuatu. (b) Benar.
(c) Salah: \omterm{def:g2:drawing-to-map:map}{peta} sama luas mempertahankan luas, tetapi mengubah bentuk,
terutama di dekat tepinya. (d) Benar.
\end{solution}

\begin{solution}{exo:g9:maps-of-the-earth:4}
(a) Lerengnya curam. (b) Tanahnya landai. (c) Sebuah \omterm{def:g1:landscapes-around-us:landform}{bukit}, dengan
puncaknya di dalam garis yang terkecil. (d) Satu titik permukaan tanah
hanya mempunyai satu \omterm{def:g4:world-physical-map:altitude}{ketinggian}, sehingga titik itu tidak dapat terletak
pada dua garis yang \omterm{def:g4:world-physical-map:altitude}{ketinggiannya} berbeda.
\end{solution}

\begin{solution}{exo:g9:maps-of-the-earth:5}
(a) Pada \ang{60}\,LU, \omterm{def:g1:landscapes-around-us:water}{danau} itu digambar 4 kali terlalu besar: \omterm{def:g1:landscapes-around-us:water}{danau} itu
tampak seluas $4 \times 1000 = \qty{4000}{km^2}$. Di \omterm{def:g3:continents-and-oceans:equator}{khatulistiwa} luas tidak
diperbesar, sehingga \omterm{def:g1:landscapes-around-us:water}{danau} yang digambar dengan ukuran yang sama mempunyai
luas sebenarnya \qty{4000}{km^2}. (b) $4 \div 1.33 \approx 3$ kali lebih
besar. (c) Luas, dan jarak dari utara ke selatan, bertambah tanpa batas ke
arah \omterm{def:g3:continents-and-oceans:poles}{kutub}: \omterm{def:g3:continents-and-oceans:poles}{kutub} akan terletak jauh tak terhingga dari \omterm{def:g3:continents-and-oceans:equator}{khatulistiwa}.
\end{solution}

\begin{solution}{exo:g9:maps-of-the-earth:6}
(a) \qty{50}{m}. (b) Puncaknya berada pada \qty{544}{m}; kontur indeks di
sekelilingnya adalah garis \qty{500}{m}. (c) Di bagian barat: \omterm{def:g1:landscapes-around-us:water}{sungai}
mengalir menuruni lereng, sehingga dasarnya lebih rendah di hilir, di
timur. (d) Di antara sekitar garis 300 dan garis \qty{400}{m}: di sana
lerengnya paling curam, dan lebih landai baik di dekat dasar \omterm{def:g1:landscapes-around-us:landform}{lembah} maupun
di dekat puncak.
\end{solution}

\begin{solution}{exo:g9:maps-of-the-earth:7}
(a) C sekitar \qty{400}{m}, D sekitar \qty{530}{m}. (b) Sekitar
\qty{135}{m}, sekitar \qty{1.3}{km} dari C, di dasar \omterm{def:g1:landscapes-around-us:landform}{lembah}. (c) Sekitar
0.3 sampai \qty{0.4}{km}. (d) \omterm{def:g4:world-physical-map:altitude}{Ketinggiannya} kecil dibandingkan dengan
jaraknya: tanpa diregangkan, \omterm{def:g1:landscapes-around-us:landform}{bukit-bukit} itu akan tampak hampir datar.
\end{solution}

\begin{solution}{exo:g9:maps-of-the-earth:8}
(a) $2 \times 25\,000 = \num{50000}$~cm $= \qty{500}{m}$. (b) $300 - 200 =
\qty{100}{m}$. (c) $100 \div 500 = 0.2$, kemiringan \qty{20}{\%}. (d) Lebih
landai: \qty{4}{cm} sama dengan \qty{1000}{m}, sehingga kemiringannya
$100 \div 1000 = \qty{10}{\%}$, setengah kecuraman tadi.
\end{solution}

\begin{solution}{exo:g9:maps-of-the-earth:9}
(a) $14 \times 18 = 252$ orbit. (b) $14 \times 103 = 1442$ menit, sekitar
\qty{24}{h}: satu hari. (c) $5.5 \div 1 = 5.5$ kali. (d) Misalnya: \omterm{def:g8:ice-on-earth:glacier}{gletser}
yang mundur, hutan yang ditebang, kota yang tumbuh, pantai atau \omterm{def:g1:landscapes-around-us:water}{sungai} yang
berpindah.
\end{solution}

\begin{solution}{exo:g9:maps-of-the-earth:10}
(a) $20\,200 \div 300\,000 \approx 0.067$~s $= 6.7 \times 10^{-2}$~s.
(b) $300\,000 \times 10^{-3} = \qty{300}{km}$. (c) Kesalahan waktu yang
sangat kecil menimbulkan kesalahan jarak yang sangat besar: jamnya harus
luar biasa teliti.
\end{solution}

\begin{solution}{exo:g9:maps-of-the-earth:11}
(a) Jarak dari titik itu ke A: $\sqrt{3.6^2 + 4.8^2} = \sqrt{12.96 + 23.04} =
\sqrt{36} = 6$; jarak ke B: $\sqrt{6.4^2 + 4.8^2} = \sqrt{40.96 +
23.04} = \sqrt{64} = 8$. (b) $(3.6, -4.8)$, bayangan cerminnya terhadap
garis AB. (c) Jarak dari C ke $(3.6, 4.8)$: $\sqrt{3.6^2 + 5.2^2} =
\sqrt{40} \approx 6.3$; ke $(3.6, -4.8)$: $\sqrt{3.6^2 + 14.8^2} =
\sqrt{232} \approx 15.2$. Penerima itu berada di $(3.6, 4.8)$. (d) Satu
jarak memberikan sebuah lingkaran (sebuah bola di ruang angkasa), dua
jarak memberikan dua titik yang mungkin; diperlukan satu jarak lagi untuk
memilih.
\end{solution}

\begin{solution}{exo:g9:maps-of-the-earth:12}
(a) $29 \div 2.1 \approx 14$ kali. (b) $2.1 \times 10.5 \approx
22$~juta~\unit{km^2}. (c) Greenland tampak meliputi sekitar 22~juta
\unit{km^2}, mendekati 29~juta milik Afrika: pada \omterm{def:g2:drawing-to-map:map}{peta} dinding Mercator,
keduanya tampak hampir sama besar, padahal Afrika 14 kali lebih luas.
(d) \omterm{def:g9:maps-of-the-earth:projection}{Proyeksi sama luas}.
\end{solution}

\begin{solution}{pb:g9:maps-of-the-earth:1}
\textbf{1.} $50\,000$~cm $= \qty{500}{m}$.

\textbf{2.} $3.6 \times 500 = \qty{1800}{m} = \qty{1.8}{km}$.

\textbf{3.} B--W: $2.5 \times 500 = \qty{1250}{m}$; W--T: $2.35 \times 500
\approx \qty{1180}{m}$.

\textbf{4.} Ke arah timur laut.

\textbf{5.} \qty{50}{m}.

\textbf{6.} S--B menurun: $296 - 143 = \qty{153}{m}$.

\textbf{7.} B--W: $335 - 143 = \qty{192}{m}$; W--T: $544 - 335 =
\qty{209}{m}$.

\textbf{8.} $209 \div 1180 \approx 0.18$, sekitar \qty{18}{\%}.

\textbf{9.} Garis lurus dari S ke T akan melintasi \omterm{def:g1:landscapes-around-us:water}{sungai} di tempat yang
tidak ada jembatannya; jalan itu turun ke jembatan di B.

\textbf{10.} $1.8 + 1.25 + 1.18 \approx \qty{4.2}{km}$ dengan kecepatan
\qty{4}{km} per jam: sekitar 63 menit; pendakian
$192 + 209 = \qty{401}{m}$: sekitar 40 menit lagi. Seluruhnya sekitar
1~jam~43~menit.

\textbf{11.} Dari T ia turun ke B, lalu mendaki dari B ke S: $296 -
143 = \qty{153}{m}$.

\textbf{12.} Panjang: $1.8 + 1.25 + 1.18 \approx \textbf{4.2 km}$; total
pendakian: $192 + 209 = \textbf{401 m}$.
\end{solution}
